\documentclass[11pt]{article}
\usepackage{amsmath,amssymb,amsthm}
\usepackage[margin=1.2in]{geometry}
\usepackage{microtype}
\usepackage{parskip}

\begin{document}

\begin{center}
  {\large\bfseries RetroDiff: Training Objective Derivation}
\end{center}

\medskip

\noindent
Let $x$ denote a clean data sample and $x_t$ its noised version at diffusion step $t \in \{1,\ldots,T\}$.
The forward process $q(x_{1:T}|x)$ gradually corrupts $x$; the model $p_\theta$ learns to reverse it.
Applying Jensen's inequality to the intractable log-likelihood yields the \emph{evidence lower bound} (ELBO):

\begin{align}
  \log p_\theta(x) \;\geq\;
  \mathbb{E}_{q(x_{1:T}|x)}\!\left[\log\frac{p_\theta(x_{0:T})}{q(x_{1:T}|x)}\right]
  = -\mathcal{L}_{\mathrm{ELBO}}
  \label{eq:elbo}
\end{align}

\noindent
Because the inequality is tight only when the variational posterior matches the true posterior,
maximising the right-hand side of \eqref{eq:elbo} is equivalent to minimising
$\mathcal{L}_{\mathrm{ELBO}}$.
Expanding the joint $p_\theta(x_{0:T})$ and applying the Markov structure of the forward process
decomposes the bound into a reconstruction term and a sum of per-step KL divergences:

\begin{align}
  \mathcal{L}_{\mathrm{ELBO}}
  = \underbrace{\mathbb{E}_q\!\left[-\log p_\theta(x_0|x_1)\right]}_{\text{reconstruction}}
  + \sum_{t=2}^{T}
    \underbrace{\mathrm{KL}\!\left(q(x_{t-1}|x_t,x) \,\|\, p_\theta(x_{t-1}|x_t)\right)}_{\text{step-}t\text{ denoising}}
  \label{eq:decomp}
\end{align}

\noindent
Each KL term in \eqref{eq:decomp} measures how well the learned reverse step $p_\theta(x_{t-1}|x_t)$
approximates the tractable posterior $q(x_{t-1}|x_t,x)$; both are Gaussian under the standard
linear noise schedule, so each term admits a closed form.

\medskip
\noindent\textbf{RetroDiff extension.}\quad
The key contribution of RetroDiff is to condition the entire objective on a set of retrieved exemplars
$R(x)$ drawn from the training corpus at inference time.
Concretely, the per-sample ELBO is evaluated with $R(x)$ appended to the model's context,
and the final training objective averages over the data distribution:

\begin{align}
  \mathcal{L}_{\mathrm{RetroDiff}}(\theta)
  = \mathbb{E}_{x \sim p_{\mathrm{data}}}\!\left[
      \mathcal{L}_{\mathrm{ELBO}}\!\left(x \,\big|\, R(x)\right)
    \right]
  \label{eq:retro}
\end{align}

\noindent
Equation~\eqref{eq:retro} reduces to the standard diffusion objective when $R(x) = \varnothing$,
making RetroDiff a strict generalisation.
The retrieval set $R(x)$ is frozen during the backward pass; only $\theta$ is updated.

\end{document}
